\contentsline {chapter}{Abstract}{iv}{chapter*.1}%
\contentsline {chapter}{Author's Declaration}{vi}{chapter*.2}%
\contentsline {chapter}{Acknowledgments}{viii}{chapter*.15}%
\contentsline {chapter}{\numberline {1}Introduction}{1}{chapter.16}%
\contentsline {section}{\numberline {1.1}Context and Motivation}{1}{section.20}%
\contentsline {section}{\numberline {1.2}Quantum Field Theory and its Classical Limit}{3}{section.23}%
\contentsline {paragraph}{Magnus expansion and the eikonal generator.}{8}{paragraph*.45}%
\contentsline {section}{\numberline {1.3}Celestial Amplitudes and Eikonalization}{12}{section.64}%
\contentsline {section}{\numberline {1.4}Computational Techniques and Methodology}{15}{section.73}%
\contentsline {section}{\numberline {1.5}Thesis in a nutshell}{22}{section.117}%
\contentsline {chapter}{\numberline {2}WEFT of Black-Hole Binaries with Dark Photon and Axion}{24}{chapter.118}%
\contentsline {section}{\numberline {2.1}Introduction}{24}{section.119}%
\contentsline {section}{\numberline {2.2}Brief review of machinery of EFTs}{28}{section.124}%
\contentsline {section}{\numberline {2.3}Review of worldline effective field theory for a point particle in GR}{30}{section.134}%
\contentsline {section}{\numberline {2.4}Conservative dynamics of binary black holes in a theory of ALP and ULV }{33}{section.147}%
\contentsline {subsection}{\numberline {2.4.1}The Gravitational bound sector}{37}{subsection.162}%
\contentsline {subsection}{\numberline {2.4.2}The Electromagnetic (EM) bound sector}{38}{subsection.169}%
\contentsline {subsection}{\numberline {2.4.3}The Proca bound sector}{39}{subsection.176}%
\contentsline {subsection}{\numberline {2.4.4}The EM-Axion bound sector}{40}{subsection.184}%
\contentsline {subsection}{\numberline {2.4.5}The pure scalar sector}{43}{subsection.194}%
\contentsline {subsection}{\numberline {2.4.6}Proca-electromagnetic interaction sector}{48}{subsection.211}%
\contentsline {section}{\numberline {2.5}The radiative dynamics}{51}{section.221}%
\contentsline {subsection}{\numberline {2.5.1}Radiative sector PN counting}{52}{subsection.229}%
\contentsline {subsection}{\numberline {2.5.2}Scalar interaction and scalar radiation}{53}{subsection.231}%
\contentsline {subsection}{\numberline {2.5.3}Electromagnetic interaction and the multipole decomposition}{60}{subsection.271}%
\contentsline {subsection}{\numberline {2.5.4}Proca interaction and multipole decomposition}{63}{subsection.283}%
\contentsline {subsection}{\numberline {2.5.5}Gravitational radiation and Multipole decomposition}{65}{subsection.293}%
\contentsline {section}{\numberline {2.6}Discussions and outlooks}{76}{section.332}%
\contentsline {section}{\numberline {2.A}Details of computation used in Section~(\ref {Sec4.2})}{77}{section.333}%
\contentsline {section}{\numberline {2.B} Details of the Feynman integral computation used in Section~(\ref {Sec4.3})}{79}{section.342}%
\contentsline {section}{\numberline {2.C} Details of computation of Feynman integral used in (\ref {Sec4}) and (\ref {Sec5})}{82}{section.359}%
\contentsline {section}{\numberline {2.D}Comments on middle vertex contribution to the radiation}{83}{section.366}%
\contentsline {paragraph}{Two-loop Fourier family.}{87}{paragraph*.387}%
\contentsline {paragraph}{\texttt {hard-hard} region.}{89}{paragraph*.405}%
\contentsline {paragraph}{\texttt {hard-soft} and \texttt {soft-hard} regions.}{89}{paragraph*.408}%
\contentsline {paragraph}{\texttt {soft-soft} region.}{90}{paragraph*.412}%
\contentsline {section}{\numberline {2.E}Useful Feynman integrals}{91}{section.425}%
\contentsline {chapter}{\numberline {3}Classical Black-Hole Scattering in Scalar-Tensor Gravity from WQFT}{93}{chapter.432}%
\contentsline {section}{\numberline {3.1}Introduction}{93}{section.433}%
\contentsline {section}{\numberline {3.2}A quick tour to worldline quantum field theory}{96}{section.436}%
\contentsline {section}{\numberline {3.3}Derivation of the Feynman rules}{98}{section.447}%
\contentsline {section}{\numberline {3.4}Computation of impulse}{101}{section.470}%
\contentsline {subsection}{\numberline {3.4.1}1PM contribution to impulse}{101}{subsection.472}%
\contentsline {subsection}{\numberline {3.4.2}2PM contribution to the impulse}{103}{subsection.488}%
\contentsline {section}{\numberline {3.5}Computation of waveform}{120}{section.568}%
\contentsline {subsection}{\numberline {3.5.1}Scalar waveform}{121}{subsection.575}%
\contentsline {subsubsection}{Scalar waveform at 1PM}{122}{subsubsection*.576}%
\contentsline {subsubsection}{Scalar waveform at 2PM}{122}{subsubsection*.582}%
\contentsline {subsection}{\numberline {3.5.2}Gravitational waveform at 2PM: contribution due to extra scalar DOF}{134}{subsection.648}%
\contentsline {section}{\numberline {3.6}Towards massive waveform}{135}{section.656}%
\contentsline {subsection}{\numberline {3.6.1}Stationary Phase (SP) approximation: the need and general procedure}{136}{subsection.663}%
\contentsline {subsection}{\numberline {3.6.2}Dealing with the massive integrals via SP approximation}{137}{subsection.675}%
\contentsline {section}{\numberline {3.7}Discussions and outlook}{143}{section.714}%
\contentsline {section}{\numberline {3.A}Sketching the derivation of the worldline action}{145}{section.715}%
\contentsline {section}{\numberline {3.B}Impulse via Eikonal: a connection to scattering amplitude}{146}{section.721}%
\contentsline {section}{\numberline {3.C}Two master integrals used for the computation of waveform}{151}{section.749}%
\contentsline {section}{\numberline {3.D}Analysis through method of regions for $\lambda _4\varphi ^4$ vertex contribution in the waveform}{155}{section.773}%
\contentsline {section}{\numberline {3.E}Massive integral coming from derivative interaction}{156}{section.778}%
\contentsline {section}{\numberline {3.F}Computation of worldline radiation diagram using large velocity approximation}{159}{section.798}%
\contentsline {section}{\numberline {3.G}Proof of the identity in \eqref {3.54kk}}{161}{section.808}%
\contentsline {section}{\numberline {3.H}Comment on the integral convergence}{162}{section.815}%
\contentsline {section}{\numberline {3.I}Comments on the algebraic properties of the Bessel curve}{166}{section.843}%
\contentsline {paragraph}{One-loop integrals, the Bessel curve, and topological recursion.}{166}{paragraph*.844}%
\contentsline {chapter}{\numberline {4}Spinning Two-Body Problem in dCS Gravity Using WQFT}{169}{chapter.865}%
\contentsline {section}{\numberline {4.1}Introduction}{169}{section.866}%
\contentsline {section}{\numberline {4.2}Chern-Simons gravity and worldline QFT}{171}{section.868}%
\contentsline {section}{\numberline {4.3}Bootstrapping Post-Minkowskian (PM) physics from Post-Newtonian (PN) data: the idea, technology and example}{176}{section.896}%
\contentsline {section}{\numberline {4.4}Computation of eikonal phase}{184}{section.937}%
\contentsline {section}{\numberline {4.5}Conclusion and outlook}{215}{section.1069}%
\contentsline {section}{\numberline {4.A}Useful Feynman integrals}{217}{section.1077}%
\contentsline {section}{\numberline {4.B}Boundary integrals in potential mode relevant for solving the differential equation}{219}{section.1088}%
\contentsline {section}{\numberline {4.C}$A(x,\epsilon )$ in $\epsilon $ factorized form}{221}{section.1099}%
\contentsline {section}{\numberline {4.D}Iterative UT integrals}{222}{section.1101}%
\contentsline {section}{\numberline {4.E}$\boldsymbol {\tilde {h}}$ coefficients}{223}{section.1103}%
\contentsline {chapter}{\numberline {5}Heterotic Footprints in Classical Gravity}{232}{chapter.1104}%
\contentsline {section}{\numberline {5.1}Introduction}{232}{section.1105}%
\contentsline {section}{\numberline {5.2}Setup and methodology}{234}{section.1106}%
\contentsline {section}{\numberline {5.3}Computation of scattering amplitude: Soft-loop expansion and classical limit}{236}{section.1120}%
\contentsline {section}{\numberline {5.4}Post-Minkowskian potential: The Lippmann-Schwinger equation and the Infrared subtraction}{249}{section.1176}%
\contentsline {paragraph}{A digression on Green's functions.}{250}{paragraph*.1183}%
\contentsline {section}{\numberline {5.5}Eikonal exponentiation and the scattering angle}{255}{section.1212}%
\contentsline {section}{\numberline {5.6}Conclusions and Discussion}{258}{section.1235}%
\contentsline {chapter}{\numberline {6}Celestial Amplitude, Shadow, and OPE in Quadratic Gravity}{260}{chapter.1236}%
\contentsline {section}{\numberline {6.1}Introduction}{260}{section.1237}%
\contentsline {section}{\numberline {6.2} Celestial eikonal amplitude for quadratic gravity}{262}{section.1238}%
\contentsline {subsection}{\numberline {6.2.1}Eikonal amplitude in GR}{266}{subsection.1257}%
\contentsline {subsection}{\numberline {6.2.2}Eikonal amplitude in quadratic gravity}{267}{subsection.1261}%
\contentsline {subsection}{\numberline {6.2.3}Analyticity and dispersion relation for celestial eikonal amplitude}{271}{subsection.1276}%
\contentsline {subsection}{\numberline {6.2.4}Dispersion Relations}{274}{subsection.1294}%
\contentsline {section}{\numberline {6.3}Shadowed correlator and operator product expansion (OPE)}{276}{section.1304}%
\contentsline {subsection}{\numberline {6.3.1}Shadowed amplitude corresponding to eikonal amplitude }{277}{subsection.1307}%
\contentsline {subsection}{\numberline {6.3.2}Shadowed amplitude corresponding to celestial Born amplitude}{278}{subsection.1312}%
\contentsline {subsection}{\numberline {6.3.3}Conformal block expansion and partial wave coefficients}{281}{subsection.1330}%
\contentsline {section}{\numberline {6.4}Connection to Carrollian amplitude}{288}{section.1359}%
\contentsline {section}{\numberline {6.5}Conclusion and Discussion}{289}{section.1364}%
\contentsline {section}{\numberline {6.A}Few definitions and conventions}{290}{section.1368}%
\contentsline {chapter}{\numberline {7}Conclusions}{293}{chapter.1378}%
\contentsline {chapter}{References}{296}{chapter.1378}%
